\documentclass[a4paper,12pt]{article}
\usepackage[UTF8]{ctex} % 支持中文
\usepackage{amsmath, amssymb, amsfonts} % 数学公式支持
\usepackage{geometry} % 页面设置
\usepackage{hyperref} % 目录超链接
\usepackage{enumitem} % 列表环境优化
\usepackage{titlesec} % 标题格式

% 页面边距设置
\geometry{left=2.5cm,right=2.5cm,top=2.5cm,bottom=2.5cm}

% 标题信息
\title{数学考研模拟题-大题集锦}
\author{}
\date{2026年2月4日}

\begin{document}

\maketitle

% 自动生成目录
\tableofcontents
\newpage

\part{第一部分 高等数学}

\section{第一章 函数、极限、连续}

\subsection{求函数的极限}

\textbf{17.(本题满分10分) 李林六套卷五}

求 $\lim_{x\to0}\frac{\int_{0}^{x}\mathrm{d}u\int_{0}^{u}\left[u^{2}-3\sin(u-t)^{2}\right]\mathrm{d}t}{x^{8}}$

\vspace{0.5cm}

\textbf{17.(本题满分10分) 余丙森五套卷一}

若函数 $y = y(x)$ 由方程 $x^{3} + y^{3} + xy - 1 = 0$ 确定，求 $\lim_{x \to 0} \frac{3y + x - 3}{x^{3}}$.

\vspace{0.5cm}

\textbf{17.(本题满分10分) 余丙森五套卷二}

设 $F(x)=\int_{0}^{x}\frac{t-\arctan t}{t}\mathrm{d}t,G(x)=\int_{0}^{x}\frac{\arcsin t-t}{t}\mathrm{d}t$，试求
\[ \lim_{x\to0}\frac{F(x)}{x^{3}},\lim_{x\to0}\frac{G(x)}{x^{3}},\lim_{x\to0}\frac{F(G(x))}{G(F(x))}. \]

\vspace{0.5cm}

\textbf{(17)（本题满分10分）合工大超越四}

已知函数 $f(x)$ 在点 $x = 0$ 处二阶可导，且 $\lim_{x \to 0} \frac{f(x) - 1 - 2x^2}{x^2} = 0$，求 $\lim_{x \to 0} \frac{f(x) - f(\sin x)}{x^4}$。

\vspace{0.5cm}

\textbf{(17)（本题满分10分）合工大超越五}

求极限 $\lim_{x\to0}\frac{(1+x)^{\frac{1}{x}}-(1+\sin x)^{\frac{1}{\sin x}}}{x-\sin x}$

\vspace{0.5cm}

\textbf{(17)(本题满分10分) 合工大超越八}

求极限 $\lim_{x\to0}\frac{\int_{0}^{x}\left[\int_{0}^{u^{2}}\left[\arctan e^{t}-\arctan(1+t)\right]dt\right]du}{\ln(1+x^{2})\int_{0}^{1}\tan(xt)^{5}dt}$

\vspace{0.5cm}

\textbf{17.(本题满分10分) 张宇八套卷二}

求极限 $\lim_{x\to0}\left[\frac{1+\int_{0}^{x}(1+t)^{\frac{1}{t}}dt}{x}-\frac{1}{\sin x}\right]$.

\vspace{0.5cm}

\textbf{17.(本题满分10分) 张宇八套卷七}

\textbf{11.(本题满分10分) 张宇八套卷七}

计算 $\lim_{n\to\infty}\sum_{i=1}^{n}\frac{1}{\sqrt{\left|2i\left(n-2i\right)\right|}}$

\subsection{已知极限反求参数}

\textbf{18.(本题满分12分) 张宇八套卷三}

已知 $\lim_{x\to+\infty}\frac{\int_{0}^{x}t^{2}e^{x^{2}-t^{2}}dt+ae^{x^{2}}}{x^{b}}=-\frac{1}{2}$，求 $a,b$ 的值.

\subsection{求数列极限(模拟卷高频压轴题)}

\textbf{21.(本题满分12分) 李林六套卷一}

设 $a_{n}=\int_{0}^{\pi}x\sin^{n}x\,dx$ （$n=1,2,\cdots$）.

（Ⅰ）证明：$a_{n}=\frac{n-1}{n}a_{n-2}$（$n=3,4,\cdots$）；

(Ⅱ) 求 $\lim_{n\to\infty}\frac{a_n}{a_{n-1}}$

\vspace{0.5cm}

\textbf{21.(本题满分12分) 李林六套卷三}

（Ⅰ）证明：方程 $\cos x = x^{n} (n = 1, 2, \cdots)$ 在 $(0, 1)$ 内有唯一实根 $x_{n}$

（Ⅱ）求 $\lim_{n\to\infty}(1-x_n)^{\frac{1}{n}\ln\cos x_n}$

\vspace{0.5cm}

\textbf{21.(本题满分12分) 李林六套卷六}

设数列 $\{x_{n}\}$ 满足 $x_{n}=\frac{1}{2}(x_{n+1}+\tan x_{n}),x_{1}\in\left(0,\frac{\pi}{4}\right),n=1,2,\cdots.$

（Ⅰ）证明：$\lim_{n\to\infty}x_{n}$ 存在，并求其值；

(Ⅱ) 求 $\lim_{n\to\infty}\left(\frac{x_{n+1}}{x_{n}}\right)^{\frac{1}{x_{n}^{2}}}$

\vspace{0.5cm}

\textbf{20.(本题满分12分) 余丙森五套卷三}

(1) 设 $0 < x < 1$，证明：$\arcsin x < \frac{\pi}{2}x$；

(2) 数列 $\{x_{n}\}$ 定义如下: $0 < x_{1} < 1, x_{n+1} = \frac{2}{\pi} \arcsin x_{n} (n = 1, 2, \cdots)$，证明: $\lim_{n \to \infty} x_{n}$ 存在并求其值.

\vspace{0.5cm}

\textbf{(19)（本题满分12分）合工大超越三}

设数列 $\{x_{n}\}$ 定义如下：$x_{1}=1, x_{n}=x_{n+1}+\mathrm{e}^{x_{n+1}}-1, (n=1,2,\cdots)$，证明数列 $\{x_{n}\}$ 的极限存在，并求出此极限值．

\vspace{0.5cm}

\textbf{(20)（本题满分12分）合工大超越五}

设 $\mathrm{d}F(x,y)=2(x-y)\mathrm{d}x-2(x+1)\mathrm{d}y$，且 $F(0,0)=4$。设 $x_{1}=4,n\geqslant1$ 时，$x_{n},x_{n+1}$ 满足 $F(x_{n},x_{n+1})=0$，证明数列 $\{x_{n}\}$ 收敛，并求极限 $\lim_{n\to\infty}x_{n}$。

\vspace{0.5cm}

\textbf{(18)（本题满分12分）}

（Ⅰ）证明方程 $\ln x + x^{n} = 0$ 在 $(0,1)$ 内有唯一实根 $x_{n}(n = 1,2,\cdots)$

（Ⅱ）证明 $\lim_{n\to\infty}x_n$ 存在.

\vspace{0.5cm}

\textbf{20.(本题满分12分) 张宇八套卷二}

设数列 $\{x_{n}\}$ 满足：$0<x_{n}<\frac{\pi}{2}$，$x_{n}\cos x_{n+1}=\sin x_{n},n=1,2,\cdots$。证明 $\{x_{n}\}$ 收敛，并求 $\lim_{n\to\infty}x_{n}$

\vspace{0.5cm}

\textbf{21. 张宇八套卷三（类似于23年选择题）}

设数列 $\{a_n\}, \{b_n\}$ 满足：$a_0 = \frac{1}{2}$, $a_{n+1} = a_n^2$, $n = 0, 1, 2, \cdots$; $\tan b_{n+1} = b_n$, $-\frac{\pi}{4} < b_n < 0$, $n = 0, 1, 2, \cdots$. 计算 $\lim_{n \to \infty} \frac{a_n}{b_n}$.

\vspace{0.5cm}

\textbf{21.(本题满分12分) 张宇八套卷四}

已知 $f(x)$ 可导，且 $\left|f'(x)\right| \leqslant \frac{1}{\mathrm{e}}$，方程 $f(x) = x$ 有唯一解 $x = 0$，又 $x_{n+1} = f(x_n)$，$n = 1，2，\cdots$。证明：当 $n \to \infty$ 时，$x_n$ 是 $e^{-\frac{n}{2}}$ 的高阶无穷小。

\vspace{0.5cm}

\textbf{18.(本题满分12分) 张宇八套卷五}

设 $S(x)=\int_{0}^{x}f(t)dt$，其中 $f(x)=\left|\arcsin(\sin x)\right|$.

(1) 写出 $f(x)$ 在 $[0, \pi]$ 上的表达式；

(2) 计算 $\lim_{x\to+\infty}\frac{S(x)}{\sqrt{1+x^2}}$

\vspace{0.5cm}

\textbf{20.(本题满分12分) 张宇八套卷六}

设正项数列 $\{x_{n}\}$ 满足 $x_{1}=\sqrt{2}, x_{n+1}^{2}=2^{x_{n}}, n=1,2,\cdots$，证明 $\{x_{n}\}$ 收敛，并求 $\lim_{n\to\infty}x_{n}$.

\vspace{0.5cm}

\textbf{20.(本题满分12分) 张宇八套卷七 (类似于 880)}

设 $a_{n}=\int_{0}^{1}(1-x^{2})^{\frac{n}{2}}dx,n=1,2,\cdots.$

(1) 证明 $a_{n+1} < a_n < a_{n-1}$;

(2) 计算 $\lim_{n\to\infty}na_{n}^{2}$

\subsection{函数的连续性及间断点类型}

\textbf{(17)(本题满分10分)}

设连续函数 $f(x)$ 满足方程 $f(x)=f\left(\frac{x}{2}\right)+\ln\cos\frac{x}{2},-\pi<x<\pi,$ 其中 $f(0)=0$.

(I) 求 $f(x)$ 的表达式；(II) 讨论 $f'(x)$ 在 $(- \pi, \pi)$ 内的连续性.

\section{第二章 一元函数微分学}

\subsection{导数定义}

\textbf{17.(本题满分10分) 张宇八套卷五}

设 $f(x)$ 在 $x=0$ 的某邻域内具有连续导数，且 $f(0)=1, f'(x)=\frac{1}{2}f[f(x)-1]$，求 $f''(0)$.

\subsection{函数的单调性、极值与最值}

\textbf{17.(本题满分10分) 李林六套卷一}

设 $f(x)$ 在 $(-\infty, +\infty)$ 内有二阶连续导数，函数 $u(x, y)$ 的全微分为 $\mathrm{d}u = y \left[ \mathrm{e}^{x} + f'(x) \right] \mathrm{d}x + f'(x) \mathrm{d}y$，且 $f(0) = f'(0) = 1$。

(Ⅰ) 求 $f(x)$;

(Ⅱ) 求 $f(x)$ 的单调区间与极值.

\vspace{0.5cm}

\textbf{17.(本题满分10分) 李林六套卷四}

设 $f(x)=\int_{0}^{\frac{\pi}{2}}|x-t|\sin t\mathrm{d}t,x\in(-\infty,+\infty)$，求 $f(x)$ 的单调区间与极值.

\vspace{0.5cm}

\textbf{(17)（本题满分10分）合工大超越二}

设 $y = y(x)$ 是由方程 $2y^{3} - 2y^{2} + 2xy - x^{2} = 1$ 确定，求 $y = y(x)$ 的极值，并求极值点处的曲率圆方程.

\vspace{0.5cm}

\textbf{(18) (本题满分12分) 合工大超越八}

设 $P(x,y)$ 是曲线 $x^{2}y-x^{2}+4xy+6y-2x-3=0$ 上的任一点，求点P在曲线最高点和最低点时的坐标.

\vspace{0.5cm}

\textbf{17.(本题满分10分) 张宇八套卷三}

设 $y = f(x)$ 由方程 $\left| x \right| y^{3} + y - 1 = 0$ 确定，求 $y = f(x)$ 的极大值.

\vspace{0.5cm}

\textbf{20.(本题满分12分) 张宇八套卷八}

设 $a_{n}=\int_{0}^{+\infty}x n^{-\frac{x}{n}}dx,n=2,3,\cdots$ 求 $\{a_{n}\}$ 的最小值.

\subsection{曲线的凹向、拐点及渐近线}

\textbf{17.(本题满分10分) 李林六套卷二}

设 $y = f(x)$ 满足 $xy' = 2y + x^{2}(x \neq 0)$，$y(e) = e^{2}$，求 $F(x) = \begin{cases} f(x), & x \neq 0, \\ 0, & x = 0 \end{cases}$ 的极值和拐点.

\vspace{0.5cm}

\textbf{18.(本题满分12分) 余丙森五套卷一}

设函数 $f(u)$ 在 $(0, +\infty)$ 内可导，$z = x f\left(\frac{y}{x}\right) + y$ 满足关系式 $x \frac{\partial z}{\partial x} - y \frac{\partial z}{\partial y} = 2z$，且 $f(1) = 1$，求曲线 $y = f(x)$ 的渐近线.

\vspace{0.5cm}

\textbf{17.(本题满分10分) 张宇八套卷一}

设 $f(x)=\lim_{n\to\infty}\frac{x}{n}\left(e^{-\frac{x^{2}}{n^{2}}}+e^{-\frac{4x^{2}}{n^{2}}}+\cdots+e^{-x^{2}}\right)$ 求：

(I) $f(x)$的表达式；(II)曲线 $y=e^{x^{2}}f(x)$的拐点.

\vspace{0.5cm}

\textbf{17 李艳芳三套卷三}

当 $x > 0$ 时, 求曲线 $y = \sqrt{x^{2} + 2x} (e^{2x} + x^{2})^{\frac{1}{x+1}}$ 的斜渐近线方程.

\subsection{证明函数不等式}

\textbf{(19)（本题满分12分）}

设 $x>0,0<\alpha<1,0<\beta<1$，证明：

（Ⅰ）函数 $f(x)=\frac{\ln(1+\alpha x)}{\ln(1+x)}$ 单调递增；

(Ⅱ) $(1 + \alpha x)(1 + \beta x) \geqslant (1 + x)^{\alpha + \beta}$.

\vspace{0.5cm}

\textbf{18.(本题满分12分) 张宇八套卷四}

设 $f(x)=x\cdot a^{x}(1-a),x>0.$

(1) 当 $0 < a < 1$ 时，求 $f(x)$ 的最大值；

(2) 当 $x > 0, 0 < y < 1$ 时，证明 $xy^{x}(1 - y) < \frac{1}{e}$.

\subsection{方程根的存在性与个数}

\textbf{(17)（本题满分10分）合工大超越三}

讨论曲线 $y = a e^{x}$ (a 为常数) 和曲线 $y = (x + 1)^{2}$ 交点的个数.

\subsection{微分中值定理有关的证明题}

\textbf{21.(本题满分12分) 李林六套卷二}

设 $f(x)$ 在 $[a,b]$ $(b>a)$ 上有二阶导数，且 $f(a)=f(b)=0, f''(x)<0.$ 证明：当 $x(a,b)$ 时，$0<f(x)<\frac{2}{b-a}\int_{a}^{b}f(x)dx$

\vspace{0.5cm}

\textbf{21.(本题满分12分) 李林六套卷四}

设 $f(x)$ 在区间 $[0,a]$ 上二阶可导，且 $\left|f''(x)\right| \leqslant M.$

（Ⅰ）若 $f(0)=f(a)=0$ ，证明：$\left|f\left(\frac{a}{2}\right)\right|\leqslant\frac{Ma^{2}}{8}$;

（Ⅱ）若 $f(x)$ 在 $(0,a)$ 内取得极值，证明：$\left|f^{\prime}(0)\right|+\left|f^{\prime}(a)\right|\leqslant Ma$

\vspace{0.5cm}

\textbf{18.(本题满分12分) 李林六套卷六}

设 $y = f(x)$ 在 $[0,1]$ 上可导，$f(x) > 0$.

（Ⅰ）证明：存在一点 $\xi \in (0,1)$，使得在区间 $[\xi,1]$ 上以 $f(\xi)$ 为高的矩形面积等于在 $[0, \xi]$ 上以 $y = f(x)$ 为曲边的梯形面积；

（Ⅱ）若 $f'(x) < 0$，证明：（Ⅰ）中的 $\xi$ 是唯一的。

\vspace{0.5cm}

\textbf{20.(本题满分12分) 余丙森五套卷一}

设函数 $f(x)$ 在 $[a,b]$ 上连续.

(1) 证明存在 $\xi \in (a, b)$，使得 $\int_{a}^{\xi} f(x) \, \mathrm{d}x = (b - \xi) f(\xi)$;

(2) 如果 $ f(x) $ 在 $(a, b)$ 内取得最大值和最小值，证明存在 $\eta \in (a, b)$，使得 $\int_{a}^{\eta} f(x) \, \mathrm{d}x = (\eta - a) f(\eta)$.

\vspace{0.5cm}

\textbf{(18)（本题满分12分）合工大超越三}

设函数 $f(x)$ 在 $[a,b]$ 上连续，且 $\int_{a}^{b}f(x)dx=0,\int_{a}^{b}xf(x)dx=m>0$ ，证明存在 $\xi\in(a,b)$，使得 $\left|f(\xi)\right|\geqslant\frac{2m}{(b-a)^{2}}$ 。

\vspace{0.5cm}

\textbf{20.(本题满分12分) 余丙森五套卷四}

设 $y=f(x)$ 在 $[0,1]$ 上非负可导，$x_{0}\in(0,1)$，且在 $[0,x_{0}]$ 上以 $f(x_{0})$ 为高的矩形面积等干函数 $f(x)$ 在 $[x_{0},1]$ 上的平均值。试证明：

(1) 存在点 $\xi \in (x_{0}, 1)$，使得 $f(\xi) = x_{0} f(x_{0})$;

(2) 对于(1)中的 $\xi$，存在 $\eta \in (0,1)$ 使得 $(\xi - x_{0})f'(\eta) = (x_{0} - 1)f(x_{0})$.

\vspace{0.5cm}

\textbf{21.(本题满分12分) 余丙森五套卷五}

(1) 设 x > 0，证明：$(x + 1)^{3} - x^{3} = 3[x + \theta(x)]^{2}$ ①，其中 $0 < \theta(x) < 1$.

(2) 证明 ① 中的 $\theta(x)$ 是 $x (0 < x < 1)$ 的减函数；

(3) 当 0 < x < 1 时，证明 ① 中的 $\theta(x)$ 满足 $\frac{1}{2} < \theta(x) < \frac{1}{\sqrt{3}}$.

\vspace{0.5cm}

\textbf{(19)(本题满分12分) 合工大超越一}

设函数$f(x)$在$[0,1]$上二阶可导,且$f(0)=0,f(1)=1$.

证明(I)存在$\xi\in(0,1)$,使$f'(\xi)=1$;$(II)$存在$\eta\in(0,1)$,使$\eta f''(\eta)+f'(\eta)=1$.

\vspace{0.5cm}

\textbf{(17) (本题满分 10 分)}

已知 $g(0)=f(0)=0, g'(0)=f'(0)=1$，且 $g''(x)-g'(x)=x-1$， $f''(x)-f'(x)\geqslant x-1$。

（Ⅰ）求函数 $g(x)$ 的表达式；（Ⅱ）证明当 x > 0 时， $f(x) \geqslant g(x)$.

\vspace{0.5cm}

\textbf{(19)（本题满分12分）合工大超越六}

设函数 $f(x)$ 在 $(-∞,+∞)$ 上三阶可导，且对任意的 $x$ 有 $\left|f(x)\right| \leqslant 1$， $\left|f''(x)\right| \leqslant 1$。

（Ⅰ）写出 $f(x+y)$ 在点 x 处的带拉格朗日余项的二阶泰勒公式；

（Ⅱ）证明对任意的 y > 0，有 $\left|f'(x)\right| \leqslant \frac{1}{y} + \frac{y^{2}}{6}$;

（Ⅲ）证明 $f'(x)$ 在 $(-∞,+∞)$ 有界，且 $\left|f'(x)\right|\leqslant\frac{\sqrt[3]{9}}{2}$

\vspace{0.5cm}

\textbf{(19)（本题满分12分）}

设 $f(t)$ 在 $[0,1]$ 上可导，且 $f'(t) > 0$.

（Ⅰ）证明对任意的 $x \in (0,1)$，存在 $\theta \in (0,1)$，使得 $\int_{0}^{x} f(t) dt = \theta x f(x) + (1 - \theta) x f(0)$;

（Ⅱ）对（Ⅰ）中的 $\theta \in (0,1)$，求 $\lim_{x\to0^{\pm}} \theta$。

\vspace{0.5cm}

\textbf{(20)（本题满分12分）合工大超越九}

设函数 $f(x)$ 在 $[0,1]$ 上二阶可导，且 $\lim_{x\to0^{+}}\frac{f(x)-f(1)}{x}=-1$，证明存在 $\xi\in(0,1)$，使得 $f''(\xi)=2$。

\vspace{0.5cm}

\textbf{(18)（本题满分12分）合工大超越十}

（Ⅰ）设 $f(x)$，$g(x)$ 均为 $[a,b]$ 上的连续函数，证明方程 $f(x)\int_{a}^{x}g(t)dt = g(x)\int_{x}^{b}f(t)dt$ 在 $(a,b)$ 内至少有一个根.

（Ⅱ）设 $f(x)$ 是 $[a,b]$ 上的可导函数且 $f'(x)>0$，证明方程 $(x-a)\int_{x}^{b}f(t)dt=(b-x)\int_{a}^{x}f(t)dt$ 在 $(a,b)$ 内没有根.

\vspace{0.5cm}

\textbf{19}

设函数 $f(x) = \begin{cases} g(x), & x > 0, \\ \int_{1}^{x^{2}} t e^{t} \, dt + \int_{x}^{-1} t^{3} e^{t^{2}} \, dt, & x \leqslant 0 \end{cases}$ 为 $(- \infty, + \infty)$ 上的连续函数，其中 $g(x)$ 为 $[0, +\infty)$ 上的可导函数，且 $g'(x) \geqslant \sqrt{1 + \cos^{2}x}$，求 $f(x)$ 在 $(- \infty, + \infty)$ 上的零点个数.

\section{第三章 一元函数积分学}

\subsection{不定积分的计算}

\textbf{17 李艳芳三套卷二}

设函数 $f(x)$ 满足 $f\left(\frac{1}{x-1}\right) = x$，且 $f(\varphi(x)) = x^3$，求 $\int\varphi(x) \, \mathrm{d}x$。

\subsection{定积分计算}

\textbf{（18）（本题满分12分）}

设 $f(x)$ 在区间 $[0,T](T>0)$ 上具有一阶连续导数， $f(0)=0,g(x)$ 在区间 $[0,T]$ 上连续.

（Ⅰ）证明 $\int_{0}^{T}f'(t)\left[\int_{t}^{T}g(x)dx\right]dt=\int_{0}^{T}f(x)g(x)dx;$

（Ⅱ）利用（Ⅰ）关系式计算积分 $I=\int_{0}^{1}\left(\int_{t}^{1}e^{-x^{2}}dx\right)dt.$

\vspace{0.5cm}

\textbf{（18）（本题满分12分）合工大超越六}

设 $f(x)$ 在 $(0, +\infty)$ 内连续， $f(1) = 1$，若对于任意的正数 a, b，积分 $\int_{a}^{ab} f(x) \, dx$ 与 a 无关，计算 $I = \int_{0}^{\frac{\pi}{2}} \frac{f(\sec^{3}x)}{\sin x + \cos x} \, dx$。

\subsection{变上限积分函数及其应用}

\textbf{21.(本题满分12分) 余丙森五套卷一}

设在区间 $(-1,+\infty)$内 $F(x)$为 $f(x)$的原函数， $f(x)F(x)=\frac{x\mathrm{e}^{x}}{2(1+x)^{2}}$，已知 $F(0)=1$，求 $f(x)$的表达式.

\vspace{0.5cm}

\textbf{21.(本题满分12分) 余丙森五套卷二}

已知 $f(x)$ 是以 $T(T>0)$ 为周期的连续函数，且 $f(0)=1, F(x)=\int_{0}^{2x} |x-t| f(t) dt$，求 $F'(T)$.

\vspace{0.5cm}

\textbf{17.(本题满分10分) 余丙森五套卷五}

设函数 $ f(x) $ 连续， $ \varphi(x) = \int_{0}^{1} f(xt) \, dt $，且 $ \lim_{x \to 0} \frac{f(x)}{x} = 1 $.

(1) 求 $\varphi'(x)$;

(2) 问 $\varphi'(x)$ 在点 x=0 处是否连续？

\vspace{0.5cm}

\textbf{(19)（本题满分12分）合工大超越十}

设 $f(x)$ 是在 $(-∞,+∞)$ 内以 $T(T>0)$ 为周期的连续函数，
\[ F(x)=\int_{0}^{x}f(t)\mathrm{d}t-\frac{x}{T}\int_{0}^{t}f(x)\mathrm{d}x \]
（Ⅰ）证明 $F(x)$ 是以 T 为周期的周期函数；（Ⅱ）计算极限 $\lim_{x\to\infty}\frac{1}{x}\int_{0}^{x}f(t)dt.$

\subsection{定积分应用}

\textbf{18.(本题满分12分) 李林六套卷三}

设容器的内表面是由曲线 $x = y + \sin y \left(0 \leqslant y \leqslant \frac{\pi}{2}\right)$ 绕 y 轴旋转一周所得的旋转曲面，以 $\pi (m^{3}/s)$ 的速率向容器内注入液体.

（Ⅰ）当液面高度为 $\frac{\pi}{4}$（m）时，求液面上升的速率；

(Ⅱ) 求将容器内注满液体需要多长时间.

\vspace{0.5cm}

\textbf{18.(本题满分12分) 李林六套卷四}

一个容器的内侧是由 $x^{2} + y^{2} = 1 \left( y \leqslant \frac{1}{2} \right)$ 绕 y 轴旋转一周而成的曲面，长度单位为 m，重力加速度为 $g \left( m/s^{2} \right)$，水的密度为 $\rho \left( kg/m^{3} \right)$，容器内盛满水.

(Ⅰ) 求容器内水的质量 m;

（Ⅱ）若将容器内的水从顶端全部抽出，至少需做功多少？

\vspace{0.5cm}

\textbf{（18）（本题满分12分）合工大超越二}

设 $A(0,0)$， $B(8,4)$是曲线 $C:y^{2}=2x$上的两个定点，P是弧AB上异于A，B的任一点，若弦PA，PB分别与曲线C所围成抛物线弓形面积为 $S_{1}$， $S_{2}$，试求点P的坐标，使 $S_{1}+S_{2}$最小，并求最小值．

\vspace{0.5cm}

\textbf{18.(本题满分12分) 李林六套卷五}

设曲线 $y=a(1-x^{2})(a>0)$ 在点 A(1,0) 和 B(-1,0) 处的法线与曲线所围封闭图形为 D.

（Ⅰ）当 D 的面积最小时，求 a 的值和最小面积；

（Ⅱ）当 D 的面积最小时，求 D 绕 y 轴旋转一周所得旋转体的体积.

\vspace{0.5cm}

\textbf{17.(本题满分10分) 余丙森五套卷四}

设函数 $ y(x) $ ($x \geqslant 0$) 由方程 $y^{3} + xy - 8 = 0$ 确定，求曲线 $ y = y(x) $ 与直线 x = 7 以及两个坐标轴所围曲边梯形绕 x 轴旋转一周所得旋转体的体积 V.

\vspace{0.5cm}

\textbf{21.(本题满分12分) 余丙森五套卷四}

某容器的外表面是 $y = x^{2} (0 \leqslant y \leqslant H)$ 绕 y 轴旋转所围成的曲面，其容积为 $450\pi \, m^{3}$，其中盛满水，如果将水全部抽出，问至少需要做多少功？

\vspace{0.5cm}

\textbf{（18）（本题满分12分）合工大超越七}

设 $f(x)=\int_{0}^{1}|t(t-x)|\mathrm{e}^{t^{2}}\mathrm{d}t$ （0<x<1），若 $f(x_{0})$ 是 $f(x)$ 的极值，求曲线 $y=f(x)$ 在点 $(x_{0},f(x_{0}))$ 处的曲率.

\vspace{0.5cm}

\textbf{19.(本题满分12分) 张宇八套卷四}

已知曲线 $L:y=\mathrm{e}^{-x}(x\geqslant0)$，设 P 是 L 上的动点，V 是 L 上从点 A(0,1) 到点 P 的一段弧绕 x 轴旋转一周所得的旋转体体积，当 P 运动到点 $\left(1,\frac{1}{\mathrm{e}}\right)$ 时，沿 x 轴正向的速度为 1，求此时 V 关于时间 t 的变化率.

\vspace{0.5cm}

\textbf{18.(本题满分12分) 张宇八套卷六}

求曲线 $y = e^{-\frac{x}{2}} \sqrt{\sin x} (x \geqslant 0)$ 绕 x 轴旋转一周所成旋转体的体积.

\vspace{0.5cm}

\textbf{21.(本题满分12分) 张宇八套卷七}

边长为 2 的等边三角形薄平板铅直沉没在水中,且一条边与水面相齐.记重力加速度为 g,水的密度为 $\rho$

(1) 求该平板一侧所受压力；

(2) 当水面开始以 0.1 的速度上涨时, 求平板一侧所受水压力的变化率.

\vspace{0.5cm}

\textbf{18}

设定义在 $(0, +\infty)$上的函数 $f(x)$满足微分方程 $y' = \frac{x^{3} + y^{3}}{xy^{2}}$，且 $f(1) = 1$.

（Ⅰ）补充定义 $f(0)$，使得 $f(x)$在 $[0,+\infty)$上连续；

（Ⅱ）求曲线 $y = [f(x)]^{3}$ 位于 x 轴下方部分的面积.

\vspace{0.5cm}

\textbf{18}

设函数 $f(x) = \frac{2x^{2} + x + 1}{x^{2} + 1} (-\infty < x < +\infty)$，记曲线 $y = f(x)$ 的水平渐近线为 L.

(Ⅰ) 求 $f(x)$ 的极值；

（Ⅱ）记曲线 $y = f(x)$ 与 L 所围成的区域为 D，求 D 绕 L 旋转一周所得旋转体的体积.

\section{第五章 多元函数的微分学}

\subsection{求多元函数的偏导数及全微分}

\textbf{18.(本题满分12分) 余丙森五套卷二}

设 $z=f\left(\mathrm{e}^{x}\cos y,\mathrm{e}^{x}\sin y\right)+g\left(xy\right)$，其中 f 具有二阶连续的偏导数，g 二阶可导，求 $\frac{\partial^{2}z}{\partial x^{2}}+\frac{\partial^{2}z}{\partial y^{2}}$.

\vspace{0.5cm}

\textbf{(20)（本题满分12分）合工大超越二}

已知 $z = xy - w, w = w(x + y, x - y)$，利用变换：$u = x + y, v = x - y$ 化简方程 $z_{xx}'' + 2z_{xy}'' + z_{yy}'' = 0$，其中涉及的二阶偏导数都连续.

\vspace{0.5cm}

\textbf{（20）(本题满分12分) 合工大超越八}

设 $z=f[\varphi(x-y),xy]+g\left(\frac{x}{y}\right)$，其中f具有二阶连续偏导数， $\varphi,g$具有二阶导数，求 $\frac{\partial^{2}z}{\partial x\partial y}$.

\subsection{求多元函数的极值}

\textbf{19.(本题满分12分) 李林六套卷一}

设 $u=u\left(\sqrt{x^{2}+y^{2}}\right)$有二阶连续偏导数，且满足 $x\frac{\partial u}{\partial x}+y\frac{\partial u}{\partial y}=2\left(x^{2}+y^{2}\right),u(0,0)=0$

(I) 求 $u(x, y)$

(II) 求 $f(x, y) = u(x, y) e^{-u(x, y)}$ 的极值.

\vspace{0.5cm}

\textbf{18.(本题满分12分) 李林六套卷二}

设 $f(x,y)=x^{3}+y^{3}-ax^{2}-by^{2}(a>0,b>0)$ 有极小值 -8，求 a,b 的值，使得 $\frac{x^{2}}{a^{2}}+\frac{y^{2}}{b^{2}}=1$ 所围面积最大.

\vspace{0.5cm}

\textbf{19.(本题满分12分) 李林六套卷三}

设 $f(x,y)$ 有二阶连续偏导数，且满足
\[ \frac{\partial^{2}f}{\partial x\partial y}=-4,f(x,0)=x^{2},\frac{\partial f}{\partial y}\bigg|_{(0,y)}=10y. \]
(Ⅰ) 求 $f(x, y)$;

（Ⅱ）求原点 $O(0,0)$ 到曲线 $f(x,y)=1$ 上的点的距离的最大值与最小值.

\vspace{0.5cm}

\textbf{19.(本题满分12分) 李林六套卷四}

设 $f(x,y)=x^{2}+y^{2}$ 在条件 $ax^{2}+2xy+by^{2}=1$ 下，于点 $\left(\frac{1}{2},-\frac{1}{2}\right)$ 处取得最大值，求 a,b 的值.

\vspace{0.5cm}

\textbf{19.(本题满分12分) 李林六套卷五}

设 $ f(x,y) $ 的全微分 $ \mathrm{d}f(x,y) = -(1 + a\mathrm{e}^{y})\sin x \, \mathrm{d}x + \mathrm{e}^{xy}(a\cos x - 1 - ay) \, \mathrm{d}y $，且 $ f(0,0) = 2 $， $ (a \neq 0) $.

(Ⅰ) 求 a 的值及 $f(x,y)$;

(Ⅱ) 求 $f(x,y)$ 的极值.

\vspace{0.5cm}

\textbf{19.(本题满分12分) 李林六套卷六}

设第一象限内的函数 $f(x,y)$ 的全微分为 $\mathrm{d}f(x,y)=(6x+2y)\mathrm{d}x+(2x+6y)\mathrm{d}y$ 且 $f(1,1)=8.$

(Ⅰ) 求 $f(x, y)$;

(Ⅱ) 在曲线 $f(x,y)=1$ 上求一点，使该点的切线与两坐标轴所围成的三角形面积最小，并求最小面积.

\vspace{0.5cm}

\textbf{（20）（本题满分12分）合工大超越三}

设 $z=z(x,y)$ 是由方程 $x^{2}-2x+y^{2}=e^{z}+z-2$ 确定的隐函数，求 $z(x,y)$ 在区域 $x^{2}+2y^{2}\leqslant4$ 内的最大值点和最小值点．

\vspace{0.5cm}

\textbf{19.(本题满分12分) 余丙森五套卷四}

已知函数 $z = z(x, y)$ 由方程 $x^{2} - xy + 2y^{2} - x - 3y + ze^{z} = 2(e^{2} - 1)$ 确定，求 $z = z(x, y)$ 的极值.

\vspace{0.5cm}

\textbf{18.(本题满分12分) 余丙森五套卷五}

设有抛物线 $L_{1}: y = x^{2} + 1$、直线 $L_{2}: y = x$.

(1) 求 $L_{1}$ 上的点到 $L_{2}$ 的最短距离 $d_{1}$;

(2) 求 $L_{1}$ 上的点和点 $P(5,0)$ 的最短距离 $d_{2}$

\vspace{0.5cm}

\textbf{(20) (本题满分 12 分)}

求 $f(x,y)=x^{3}-x^{2}y+\frac{5}{6}y^{2}+\frac{1}{3}y^{3}$ 的极值.

\vspace{0.5cm}

\textbf{（20）（本题满分12分）合工大超越四}

设 $P(x)=x^{2}+ax+b$ ，问 $\int_{-1}^{1}P^{2}(x)dx$ 是否有最大、最小值，若有，则求出最值，若没有，请说出理由.

\vspace{0.5cm}

\textbf{(19)(本题满分12分) 合工大超越五}

曲线 $x^{\overline{2}}+y^{\overline{2}}=1(x>0,y>0)$上任意一点处的切线在两坐标轴上的截距之和为l，试求l的最小值.

\vspace{0.5cm}

\textbf{(20)（本题满分12分）合工大超越六}

设 $f(x,y)=x^{2}-4xy+2y^{2}$

（Ⅰ）证明 $f(x,y)$ 在任意开区域内不存在最值；

（Ⅱ）求函数在条件 $x^{2}+2y^{2}=1$ 上的最大值和最小值.

\vspace{0.5cm}

\textbf{(20)（本题满分12分）合工大超越七}

设 $x^{2}-xy+y^{2}-3x+(z-1)e^{z}+3=0$ 确定了 $z=z(x,y)$，求 $z(x,y)$ 的极值.

\vspace{0.5cm}

\textbf{（20）(本题满分12分) 合工大超越十}

求 $f(x,y,z)=\frac{1}{x^{2}}+\frac{1}{y^{2}}+\frac{16}{z^{2}}$在条件 $x^{2}+y^{2}+\frac{z^{2}}{4}=1(x>0,y>0,z>0)$下的最小值.

\vspace{0.5cm}

\textbf{18.(本题满分12分) 张宇八套卷二}

已知 $f(x,y)$ 满足 $\mathrm{e}^{-2x}\frac{\partial f}{\partial x}=4y+2y^{2}+2x+1$ ，且 $f(0,y)=2y+y^{2}$ 。求：

(1) $f(x,y)$的表达式；

(2) $f(x,y)$的极值.

\vspace{0.5cm}

\textbf{19.(本题满分12分) 张宇八套卷三}

求函数 $u = \sqrt{x^{2} + y^{2} + z^{2}}$ 在约束条件 $x + 2y = 1$ 与 $x^{2} + 2y^{2} + z^{2} = 1$ 下的最值.

\vspace{0.5cm}

\textbf{19.(本题满分12分) 张宇八套卷五}

设函数 $f(x,y)=2\mathrm{e}^{x^{2}y}-\mathrm{e}^{x}-\mathrm{e}^{-x}$

(1) 计算 $\lim_{x\to0}\frac{f(x,y)}{x^{2}}$;

(2) $f(x,y)$ 在点 $(0,0)$ 处是否取得极值？若是，求出此极值，若不是，说明理由.

\vspace{0.5cm}

\textbf{18.(本题满分12分) 张宇八套卷七}

求函数 $f(x,y)=x^{3}+y^{2}+6xy$ 的极值.

\vspace{0.5cm}

\textbf{18.(本题满分12分) 张宇八套卷八}

设函数 $u = xz + ay^{3} (z \geqslant 0)$，且 $x^{2} + y^{2} + z^{2} = 1$.

(1) 当 $a = \frac{1}{3}$ 时，求 u 的最大值；

(2) 当 $a = t$ (t 为变量) 时，u 是否有最大值，若有，求出最大值，若没有，说明理由.

\vspace{0.5cm}

\textbf{20}

已知函数 $ f(u, v) $ 可微, $ g(x, y) = f(x, x + y) $, $ \frac{\partial g(x, y)}{\partial x} = x^5 + (x + y)^5 - 8x - 4y $, $ \frac{\partial g(x, y)}{\partial y} = (x + y)^5 - 4x - 2y $, 且 $ f(0, 0) = 1 $.

(Ⅰ) 求 $f(u, v)$;

(Ⅱ) 求 $f(u,v)$ 的极值.

\vspace{0.5cm}

\textbf{20 李艳芳三套卷三}

求函数 $u = x + y + z$ 在条件 $xy + z = 1$ 以及 $x^2 + y^2 + \frac{z^2}{2} = 8$ 下的极值.

\subsection{反问题}

\textbf{18.(本题满分12分) 余丙森五套卷三}

设二元函数 $z = z(x, y)$ 的全微分为
\[ \mathrm{d}z=(2xy^{3}+a\mathrm{e}^{y}\sin x)\mathrm{d}x+(bx^{2}y^{2}+\mathrm{e}^{y}\cos x)\mathrm{d}y, \]
且 $z(0,0)=1$ 。求 $z=z(x,y)$ 的表达式.

\vspace{0.5cm}

\textbf{（19）（本题满分12分）合工大超越九}

设 $z = u(x, y) e^{ax + by}$，其中 $u(x, y)$ 满足 $\frac{\partial^{2}u}{\partial x \partial y} = 0$，试确定 a, b 使函数 z 满足方程 $\frac{\partial^{2}z}{\partial x \partial y} - \frac{\partial z}{\partial x} - \frac{\partial z}{\partial y} + z = 0.$

\vspace{0.5cm}

\textbf{20 李艳芳三套卷二}

已知 $f(t)$ 为具有二阶连续导数的一元函数，且 $f(1) = 2$。若函数 $g(x, y) = x f\left(\frac{y}{x}\right) + \frac{1}{2} f\left(\frac{x}{y}\right)$，且满足 $y^2 \frac{\partial^2 g}{\partial y^2} - x^2 \frac{\partial^2 g}{\partial x^2} = \frac{x}{y} f\left(\frac{x}{y}\right) - \frac{x^2}{y^2}$，求 $g(x, y)$。

\section{第六章 重积分}

\subsection{二重积分的基本计算}

\textbf{(21) (本题满分 12 分) 合工大超越二}

(I) 设 $\mathrm{d}f(x,y)=\frac{y(1+y^{2})\mathrm{d}x+x\mathrm{d}y}{\sqrt{(1+y^{2})^{3}}}$，且 $f(0,0)=0$，求 $f(x,y)$;

(Ⅱ)计算二重积分 $I=\iint_{D}f(x,y)d\sigma$，其中区域 D 由 $y=\sqrt{x}$， $y=2\sqrt{x}$ 和 x=1 所围成.

\vspace{0.5cm}

\textbf{（21）(本题满分12分) 合工大超越四}

设 $f(x)$ 具有连续导数，且满足 $f(1)=1$ 。区域 D 是由 x 轴，y 轴和 $x+y=1$ 所围区域，若 $f(x+y)=\frac{x}{1+x+y}+\iint_{\Omega}y f'(x+y)d\sigma$ ，求 $\iint_{\Omega}f(x+y)d\sigma$ 。

\subsection{利用积分区域对称性或函数奇偶性计算积分}

\textbf{20.(本题满分12分) 李林六套卷三}

设 $D=\left\{(x,y)\mid1\leqslant x+y\leqslant2,x\geqslant0,y\geqslant0\right\}$，计算：$I=\iint_{D}\frac{x\mathrm{e}^{(x+y)^{2}}}{x+y}\mathrm{d}x\mathrm{d}y.$

\vspace{0.5cm}

\textbf{19.(本题满分12分) 余丙森五套卷三}

设 D 是由直线 $x + y = 1, x + y = 2$ 及 x 轴和 y 轴围成的四边形区域，计算二重积分 $I = \iint_{D} e^{(x + y)^{2}} (\cos^{2} x + \sin^{2} y) \, dx \, dy.$

\vspace{0.5cm}

\textbf{19.(本题满分12分) 余丙森五套卷五}

设函数 $f(x)$ 连续，
\[ f(x)=x^{2}+x\int_{0}^{x^{2}}f(x^{2}-t)\mathrm{d}t+\iint_{D}f(xy)\mathrm{d}x\mathrm{d}y, \]
其中 D 为以 $A(1,-1)$, $B(-1,1)$, $C(1,1)$ 为顶点的三角区域，若 $f(1)=0$，

(1) 求 $\int_{0}^{1} f(x) \, dx + \iint_{D} f(xy) \, dx \, dy$ 的值；

(2) 求 $\int_{0}^{1} f(x) \, dx$ 的值.

\vspace{0.5cm}

\textbf{(21)(本题满分12分) 合工大超越七}

设D是由直线 $y=\sqrt{3}x$, $x=\sqrt{3}y$及曲线 $xy=1$, $xy=2$围成的位于第一象限的四边形区域.(I)证明: $\arctan x+\arctan\frac{1}{x}=\frac{\pi}{2}$，其中x>0；(II)计算二重积分 $I=\iint_{D}e^{xy}\arctan\frac{y}{x}d\sigma$

\vspace{0.5cm}

\textbf{(21)（本题满分12分）合工大超越九}

计算积分 $I=\iint_{D}\frac{1}{\left(a^{2}+x^{2}+y^{2}\right)^{\frac{3}{2}}}\mathrm{d}\sigma$ ，其中 $D:0\leqslant x\leqslant a,0\leqslant y\leqslant a.$

\vspace{0.5cm}

\textbf{(21) (本题满分 12 分) 合工大超越十}

设 $D = \{(x, y) \mid 1 \leqslant |x| + |y| \leqslant 2\}$，计算 $I = \iint_{D} \frac{e^{(|x| + |y|)^2} \cos x}{\cos x + \cos y} \, \mathrm{d}\sigma.$

\vspace{0.5cm}

\textbf{20.(本题满分12分) 张宇八套卷一}

设函数 $f(x)$ 满足 $f(x)=x\int_{0}^{x^{2}}f(x^{2}-t)dt+\iint_{D}f(x+y)d\sigma+a$，其中 D 是由 $y=x^{3}$ 与 y=1，y=-1 及 y 轴所围平面有界闭区域， $f(1)=0$，且 $f(x)$ 在 $[0,1]$ 上的平均值为 3，求常数 a 的值.

\vspace{0.5cm}

\textbf{20.(本题满分12分) 张宇八套卷五}

计算 $\iint_{D}\max\{x,y\}d\sigma$，其中D是 $x^{2}+y^{2}\leqslant2x$与 $x^{2}+y^{2}\leqslant2y$重合的部分.

\vspace{0.5cm}

\textbf{19 李艳芳三套卷一}

设平面区域由 $x^{\frac{2}{3}} + y^{\frac{2}{3}} = 1 (x \geqslant 0, y \geqslant 0)$ 与两坐标轴围成，计算 $\iint_{D} (\sqrt[3]{x} + \sqrt[3]{y}) \, dx \, dy$.

解 曲线 $x^{\frac{2}{3}} + y^{\frac{2}{3}} = 1 (x \geqslant 0, y \geqslant 0)$ 为星形线位于第一象限的部分，其参数方程为 $\left\{\begin{aligned} x &= \cos^{3}t, \\ y &= \sin^{3}t \end{aligned}\right.$ $\left(0 \leqslant t \leqslant \frac{\pi}{2}\right)$.

\subsection{分区域积分的二重积分}

\textbf{19.(本题满分12分) 李林六套卷二}

设 $f(x,y)=\left\{\begin{aligned}&x-y+1,&x+y\leqslant1\\ &\frac{(x+y)^{2}}{x^{2}+y^{2}},&x+y>1\end{aligned}\right.$， $D=\left\{(x,y)\mid x+y+xy\geqslant1,x^{2}+y^{2}\leqslant1\right.$， $x\geqslant0,y\geqslant0$，计算 $I=\iint_{D}f(x,y)dxdy.$

\vspace{0.5cm}

\textbf{20.(本题满分12分 李林六套卷四}

设 $D=\left\{(x,y)\mid x^{2}+y^{2}\leqslant4\right\}$，计算 $\iint_{D}|2x-x^{2}-y^{2}|\,\mathrm{d}x\,\mathrm{d}y.$

\vspace{0.5cm}

\textbf{19.(本题满分12分) 余丙森五套卷一}

计算二重积分 $I=\iint_{D}f(x,y)\mathrm{d}x\mathrm{d}y,\quad D:0\leqslant x\leqslant2,\left|y\right|\leqslant x.$ 其中
\[ f(x,y)=\min_{t\geqslant1}(t^{4}-4x^{3}t+y^{3}). \]

\vspace{0.5cm}

\textbf{(21)（本题满分12分）合工大超越一}

设 $f(x,y)=\left\{\begin{aligned}&\sqrt{\frac{1-x^{2}-y^{2}}{1+x^{2}+y^{2}}},&y\leqslant x,\\&x^{2}y,&y>x,\end{aligned}\right.$ $D:x^{2}+y^{2}\leqslant1,x\geqslant0$ ，计算积分 $\iint_{D}f(x,y)d\sigma.$

\vspace{0.5cm}

\textbf{(21)（本题满分12分）合工大超越六}

计算：$I=\iint_{D}\left|x+y\right|-1\left|\mathrm{d}x\mathrm{d}y\right.$，其中 $D=\left\{(x,y)\mid x^{2}+y^{2}\leqslant1,0\leqslant x\leqslant1\right\}$

\vspace{0.5cm}

\textbf{20.(本题满分12分) 张宇八套卷三}

设区域 D 为 $x^{2} + y^{2} \leqslant 2x$ 与 $|y| \leqslant \sqrt{3}x$ 的重合部分，计算
\[ \iint_{D}\max\{\sqrt{2-(x^{2}+y^{2})},\sqrt{x^{2}+y^{2}}\}\mathrm{d}x\mathrm{d}y. \]

\vspace{0.5cm}

\textbf{19.(本题满分12分) 张宇八套卷八}

计算二重积分 $\iint_{D}\left|x^{2}+y^{2}-\sqrt{2}(x+y)\right|\mathrm{d}x\mathrm{d}y$，其中 $D=\{(x,y)\mid x^{2}+y^{2}\leqslant4\}$

\subsection{交换积分次序及坐标系}

\textbf{20.(本题满分12分) 李林六套卷一}

设D是由曲线 $y=\sqrt{1-x^{2}}$和 $y=\sqrt{4-x^{2}}$与 $y\geq-x$及x轴所围成的区域，

计算 $I=\iint\limits_{D}\frac{\sqrt{x^{2}+y^{2}}}{x^{2}+2y^{2}}dx dy.$

\vspace{0.5cm}

\textbf{21.(本题满分12分) 李林六套卷五}

设 $f(x)$ 在 $[0,1]$ 上有连续导数， $f(0)=1$， $f_{+}^{\prime}(0)=1$。且
\[ \iint_{D_{t}}f^{\prime\prime}(x+y)\mathrm{d}x\mathrm{d}y=\iint_{D_{t}}\left[f(x+y)+(x-y)\right]\mathrm{d}x\mathrm{d}y. \]
其中 $D_{t}=\left\{(x,y)\mid0\leqslant y\leqslant t-x,0\leqslant x\leqslant t\right\}(0<t\leqslant1)$，求 $f(x)$.

\vspace{0.5cm}

\textbf{20.(本题满分12分) 李林六套卷六}

设 $D=\left\{(x,y)\mid x^{2}+y^{2}\leqslant1,0\leqslant y\leqslant x\right\}$，计算 $I=\iint_{D}\frac{xy}{1+x^{2}-y^{2}}dx dy.$

\vspace{0.5cm}

\textbf{(21)（本题满分12分）合工大超越三}

（Ⅰ）当x>0时，求极限 $\lim_{x\to0}\frac{x^{2}+y^{2}}{\sqrt{x^{2}+y^{2}}+x}$.

（Ⅱ）计算二重积分 $I=\iint_{D}\frac{x^{2}+y^{2}}{\sqrt{x^{2}+y^{2}}+x}d\sigma$，其中 D 是由 $x=\sqrt{y-y^{2}},y=\sqrt{1-x^{2}}$ 以及 x 轴所围位于第一象限区域.

\vspace{0.5cm}

\textbf{(21)（本题满分12分）合工大超越五}

设 $D=\left\{(x,y)\mid4x^{2}y^{2}\leqslant(x^{2}+y^{2})^{3}\leqslant16x^{2}y^{2},x\geqslant0,y\geqslant0\right\}$，计算积分 $I=\iint_{D}(x^{2}+y^{2}+xy)dxdy$.

\vspace{0.5cm}

\textbf{(21)（本题满分12分）合工大超越八}

利用二重积分计算定积分：$I=\int_{0}^{1}x^{2}f(x)dx$，其中 $f(x)=\int_{x^{3}}^{x}e^{-y^{2}}dy.$

\vspace{0.5cm}

\textbf{21.(本题满分12分) 张宇八套卷二}

设平面区域 $D=\left\{(x,y)\mid x+y\leqslant1,x\geqslant0,y\geqslant0\right\}$，计算 $\iint_{D}\frac{e^{-(x+y)}}{\sqrt{xy}}d\sigma$

\vspace{0.5cm}

\textbf{20.(本题满分12分) 张宇八套卷四}

设平面区域 $D=\left\{(x,y)\mid(x-1)^{2}+y^{2}\geqslant1,(x-2)^{2}+y^{2}\leqslant4,y\geqslant x\right\}$，计算 $\iint_{D}(x^{2}+y^{2})d\sigma$

\vspace{0.5cm}

\textbf{18 李艳芳三套卷三}

设平面有界区域 D 位于第一象限，由曲线 $(x^{2} + y^{2})^{4} = x^{6} + y^{6}$ 与 $4(x^{2} + y^{2})^{4} = x^{6} + y^{6}$ 以及坐标轴围成，计算 $\iint_{D} \frac{\sqrt{x^{6} + y^{6}}}{(x^{2} + y^{2})^{2}} \, \mathrm{d}x \, \mathrm{d}y$。

\subsection{二重积分证明}

\textbf{19. 张宇八套卷六}

设平面区域 $D = \left\{(x, y) \mid 2x^2 + y^2 \leqslant 2\sqrt{x^2 + y^2}, y \geqslant x \geqslant 0\right\}$，若 $\iint_D \frac{f(x, y)}{\sqrt{1 - x^2}} \, \mathrm{d}\sigma = a > 0$， $f(x, y)$ 是 $D$ 上的连续函数。

(1) 计算 $\iint_D \frac{1}{\sqrt{1 - x^2}} \, \mathrm{d}\sigma$。

(2) 证明: 存在 $(\xi, \eta) \in D$，使得 $|f(\xi, \eta)| \geq \frac{\sqrt{2}}{\pi} a$.

\vspace{0.5cm}

\textbf{21 李艳芳三套卷二}

设函数 $ f(x, y) $ 具有一阶连续偏导数，且 $ f(0, 0) = 0 $， $ |f_1'(x, y)| \leqslant 4 |x - y| $， $ |f_2'(x, y)| \leqslant 4 |x - y| $，区域 D 为由直线 y = x，x = 1 以及 x 轴所围成的平面区域。证明：$ \left| \iint_{D} f(x, y) \, \mathrm{d}x \, \mathrm{d}y \right| \leqslant \frac{1}{6} $。

\section{第九章 常微分方程}

\subsection{一阶微分方程}

\textbf{17.(本题满分10分) 李林六套卷六}

设 $y = y(x)$ 满足 $x^{2}y' + y = x^{2}e^{\frac{1}{x}}(x \neq 0)$，且 $y(1) = 3e$.

(Ⅰ) 求 $y = y(x)$ 的全部渐近线方程；

（Ⅱ）讨论曲线 $y = y(x)$ 与 $y = k(k > 0)$ 不同交点的个数.

\vspace{0.5cm}

\textbf{18.(本题满分12分) 余丙森五套卷一}

设函数 $ f(u) $ 在 $(0, +\infty)$ 内可导， $ z = xf\left(\frac{y}{x}\right) + y $ 满足关系式 $x \frac{\partial z}{\partial x} - y \frac{\partial z}{\partial y} = 2z$，且 $ f(1) = 1 $，求曲线 $ y = f(x) $ 的渐近线.

\vspace{0.5cm}

\textbf{(17)（本题满分10分）}

设 $f(x)$ 在 $(-\infty, +\infty)$ 内有定义， $f'(0) = 1$，且对任意 $x, y \in (-\infty, +\infty)$ 恒有 $f(x + y) = \mathrm{e}^{y} f(x) + \mathrm{e}^{x} f(y)$.

（Ⅰ）求 $f(x)$ 满足的一阶微分方程；（Ⅱ）求 $f(x)$ 表达式.

\vspace{0.5cm}

\textbf{19.(本题满分12分) 张宇八套卷一}

已知微分方程 $e^{y}=t+\frac{1}{y}$ 满足 $y(0)=0.$

(1) 求该微分方程的特解 $y = y(t)$;

(2) 设 $\left\{\begin{aligned} x &= \sqrt{1 + t^{2}}, \\ y &= y(t), \end{aligned}\right.$ 计算 $\left.\frac{d^{2}y}{dx^{2}}\right|_{\iota=1}$.

\vspace{0.5cm}

\textbf{21.(本题满分12分) 张宇八套卷八}

已知函数 $y = y(x)$ 满足
\[ x(\ln x-1)y^{\prime}(x)+(3-\ln x^{2})y(x)=0,x>e, \]
且 $y(\mathrm{e}^{2})=\frac{\mathrm{e}^{4}}{2}$，求 $y=y(x)$ 的最小值.

\vspace{0.5cm}

\textbf{18}

设定义在 $(0, +\infty)$上的函数 $f(x)$满足微分方程 $y' = \frac{x^3 + y^3}{xy^2}$，且 $f(1) = 1$.

（Ⅰ）补充定义 $f(0)$，使得 $f(x)$在 $[0,+\infty)$上连续；

（Ⅱ）求曲线 $y = [f(x)]^{3}$ 位于 x 轴下方部分的面积.

\subsection{高阶常系数微分方程的求解}

\textbf{17.(本题满分10分) 李林六套卷一}

设 $f(x)$ 在 $(-∞, +∞)$ 邻开六套卷一 函数 $u(x, y)$ 的全微分为 $\mathrm{d}u = y \left[ \mathrm{e}^{x} + f'(x) \right] \mathrm{d}x + f'(x) \mathrm{d}y$，且 $f(0) = f'(0) = 1$.

(Ⅰ) 求 $f(x)$;

(Ⅱ) 求 $f(x)$ 的单调区间与极值.

\vspace{0.5cm}

\textbf{20.(本题满分12分) 李林六套卷二}

设 $y = y(x)$ 由参数方程 $\left\{\begin{aligned} x &= t^{2} \\ y &= y(t) \end{aligned}\right.$ $(t > 0)$ 确定，其中 $y(x)$ 有二阶导数，且满足 $4x \frac{\mathrm{d}^{2}y}{\mathrm{d}x^{2}} + 2(1 - \sqrt{x}) \frac{\mathrm{d}y}{\mathrm{d}x} - 6y = \mathrm{e}^{-\sqrt{x}}$，若 $\lim_{x \to 0^{+}} y(x) = 0$， $\lim_{x \to +\infty} y(x) = 0$，求 $y = y(x)$.

由 $\lim_{x\to0^{+}}y(x)=0$ , 得 $C_{1}+C_{2}-\frac{1}{4}=0$ 。由 $\lim_{x\to+\infty}y(x)=0$ ，得 $C_{2}=0$ ，故 $C_{1}=\frac{1}{4}, C_{2}=0$ 。

所求函数 $y(x)=\frac{1}{4}e^{-2\sqrt{x}}-\frac{1}{4}e^{-\sqrt{x}}$

\vspace{0.5cm}

\textbf{18.(本题满分12分) 余丙森五套卷四}

设函数 $f(x)$ 在 $(-∞,+∞)$ 上连续，且满足 $f(x)=\sin x+\int_{0}^{x}tf(x-t)dt$

(1) 求 $f(x)$;

(2) 若 $xy' + y = f(x)$，求 $y = y(x)$ 的表达式.

\vspace{0.5cm}

\textbf{20.(本题满分12分) 余丙森五套卷五}

利用变换 $u = \tan y$ 将方程
\[ y^{\prime \prime}\sec^{2}y+2(y^{\prime})^{2}\sec^{2}y\tan y+2y^{\prime}\sec^{2}y-3\tan y=-3x+2 \]
化简，并求出原方程的通解.

\vspace{0.5cm}

\textbf{（18）（本题满分12分）合工大超越一}

函数 $y(x)$ 是方程 $y'' - 2y' + y = f(x)$ 满足 $y(0) = 1, y'(0) = 2$ 的特解，其中 $f(x) = \left\{ \begin{aligned} & x + 1, & x \leqslant 0, \\ & e^{x}, & x > 0. \end{aligned} \right.$

(Ⅰ) 求 $y(x)$ 的表达式；（Ⅱ）求曲线 $y(x)$ 的渐近线.

\vspace{0.5cm}

\textbf{（19）（本题满分12分）合工大超越四}

设函数 $y = f(t)$ 有二阶连续导数且 $f(t) \neq 0$，若 $f(0) = 2, f'(0) = 1$， $x = \int_{0}^{t} f(u) \, du$，且 $y^{2} \frac{d^{2}y}{dx^{2}} = \frac{dy}{dx}$，求 $f(t)$ 的表达式.

\vspace{0.5cm}

\textbf{(17)（本题满分10分）合工大超越九}

设函数 $y = y(x)$ 由参数方程 $\left\{\begin{aligned}y &= \varphi(t)\\ y &= \varphi(t)\end{aligned}\right.$ 确定，其中 $\varphi(t)$ 具有二阶导数且 $y = \varphi(t)$ 在点 $(0,1)$ 处有水平切线. 若 $\frac{d^{2}y}{dx^{2}} = 4t^{2} + 6t + 2$，试求 $y = \varphi(t)$ 的表达式.

\vspace{0.5cm}

\textbf{21.(本题满分12分) 张宇八套卷六}

若函数 $f(x)$ 满足关系式 $f'(x) + af(x) = \int_{x}^{0} f(t) \, dt, a > 0$，求 $\int_{0}^{+\infty} f(x) \, dx$.

\subsection{微分方程的几何/物理应用}

\textbf{18.(本题满分12分) 李林六套卷一}

经过曲线 $y = y(x)$ 上点 $P(x, y)$ 作 x 轴垂线交于点 $M(x, 0)$，曲线上点 $P(x, y)$ 处的切线为 PT，已知点 M 到切线 PT 的距离等于 1，且 $y(0) = 1$.

(Ⅰ) 求 $y = y(x)$;

（Ⅱ）求曲线 $y = y(x)$， $x \in [-1, 1]$ 的弧长.

\vspace{0.5cm}

\textbf{17.(本题满分10分) 李林六套卷三}

设 $f(x)$ 在 $(0, +\infty)$ 内有定义， $f'(1) = 1$，且对任意的 $x, y \in (0, +\infty)$，有 $f(xy) = f(x) + f(y)$.

（Ⅰ）证明：$f(x)$ 在 $(0,+\infty)$ 内可导，并求 $f(x)$;

（Ⅱ）求曲线 $y = f(x)$ 的一条切线，使该切线与 x = 2, x = 6 及 $y = f(x)$ 所围成的图形的面积最小.

\vspace{0.5cm}

\textbf{20.(本题满分12分) 李林六套卷五}

设 $y=f(x)$ 在 $[0,+\infty)$ 上可导， $f(0)=k>0$，当 x>0 时， $f'(x)>0$，若在区间 $[0,x]$ 上以 $y=f(x)$ 为曲边的曲边梯形面积为 $A(x)$， $y=f(x)$ 在 $[0,x]$ 上的弧长为 $S(x)$，且 $A(x)=kS(x)$，求 $f(x)$.

\vspace{0.5cm}

\textbf{21.(本题满分12分) 余丙森五套卷三}

设 C 为上半平面内的一条凹的曲线，已知 C 过点 $P_{0}(0,1)$，其上任一点 $P(x,y)$ 处的法线与 x 轴的交点为 Q，又该曲线在 $P(x,y)$ 处的曲率等于此曲线在该点处的法线上线段 PQ 长度的倒数，且曲线在点 $P_{0}$ 处的切线与 x 轴平行，求此曲线方程.

\vspace{0.5cm}

\textbf{(18)(本题满分12分) 合工大超越五}

设曲线 $y=y(x)(x\geq0)$是第一象限的一条凸曲线,曲线上点 $(0,1)$处的切线方程为 $y=x+1$，且其上任一点 $P(x,y)$处的切线和y轴以及曲线 $y(x)$围成图形的面积恒等于 $\int_{0}^{x}t^{2}[y'(t)+1]dt$.求该曲线的方程.

\vspace{0.5cm}

\textbf{19.(本题满分11分) 张宇八套卷二}

已知函数 $y = y(x)$ 可导，将区间 $[0, x]$ 上以 $y = y(x)$ 为曲边的曲边梯形面积记为 $S$，在 $[0, x]$ 的弧长记为 $aS(a > 0)$，且 $y(0) = \frac{1}{a}$，求曲线 $y = y(x)$ 的表达式。

\vspace{0.5cm}

\textbf{19 李艳芳三套卷三}

已知函数 $f(x)$ 为定义在 $[0, +\infty)$ 上的非负可导函数， $f'(1) = 8(e^{-1} + 1)$. 设点 $(X, Y)$ 为曲线 $y = f(x)$ 上任意一点，记曲线在该点处的切线在 y 轴上的截距为 A，该曲线与直线 x = X 以及坐标轴所围区域绕 y 轴旋转一周所得旋转体的体积为 B. 若 $A + \frac{B}{2\pi} = X^{4}$，求 $f(x)$.

\part{第二部分 线性代数}

\section{第一类 对角化、化标准形}

\textbf{22.(本题满分12分) 余丙森五套卷三}

已知二次型 $f(x_{1},x_{2},x_{3})=(x_{1},x_{2},x_{3})\begin{pmatrix}1-a&2a&-1\\2&1-a&1\\1&-1&2\end{pmatrix}\begin{pmatrix}x_{1}\\x_{2}\\x_{3}\end{pmatrix}$ 的秩为 2，

(1) 求常数 a 的值;

(2) 求一个正交变换 x = Qy 化二次型为标准形；

(3) 求方程 $f(x_{1}, x_{2}, x_{3}) = 0$ 的解.

\vspace{0.5cm}

\textbf{（22）（本题满分12分）合工大超越二}

设3阶实对称矩阵A有重特征值 $\lambda(\lambda\neq0)$，且A的迹 $\mathrm{tr}(\mathbf{A})=2$，矩阵 $A^{T}$的列向量均与列向量 $(1,1,1)^{\mathrm{T}}$正交.

（Ⅰ）求矩阵 A 的所有特征值及特征向量；

（Ⅱ）求正交阵 Q 将 A 对角化；

（Ⅲ）二次型 $f(x_{1},x_{2},x_{3})=\boldsymbol{x}^{\mathrm{T}}\boldsymbol{A}\boldsymbol{x}$，求 f=0 的解.

\vspace{0.5cm}

\textbf{(22)（本题满分12分）合工大超越三}

设A为3阶实对称矩阵， $\boldsymbol{\alpha}=(1,1,0)^{\mathrm{T}},\boldsymbol{\beta}=(1,0,1)^{\mathrm{T}},A\boldsymbol{\alpha}=2\boldsymbol{\beta},A\boldsymbol{\beta}=2\boldsymbol{\alpha}$，且存在3阶矩阵B，有 $\mathrm{r}(AB)<\mathrm{r}(B)$.

（Ⅰ）求正交变换 x=Qy，化二次型 $x^{T}Ax$ 为标准形，并写出标准形；

（Ⅱ）若 $\gamma=(1,0,4)^{\mathrm{T}}$，求 $A^{2024}\gamma$。

\vspace{0.5cm}

\textbf{(22)（本题满分12分）合工大超越四}

设A,B均为3阶方阵，其中A为实对称矩阵， $A=(a_{ij})_{3\times3},a_{11}=2,\quad B=\begin{pmatrix}1&0&-1\\-1&0&1\\0&2&0\end{pmatrix}$，且满足 $(A-E)B=O.$

（Ⅰ）求矩阵 A；（Ⅱ）求可逆矩阵 P，使 $P^{-1}AP=\Lambda$。

\vspace{0.5cm}

\textbf{(22)（本题满分12分）合工大超越六}

记 $A=\begin{pmatrix}1&2&0\\2&2&a\\0&a&3\end{pmatrix}$，设 $\xi_{1}=\begin{pmatrix}-2\\2\\1\end{pmatrix}$是方阵A的特征向量.

（Ⅰ）求 a；（Ⅱ）求方阵 B，使得 $A = B^{3}$

\vspace{0.5cm}

\textbf{(22)（本题满分12分）合工大超越八}

已知 $A=\begin{pmatrix}-1&a&0\\0&b&2\\0&0&3\end{pmatrix},B=\begin{pmatrix}2&-6&1\\0&-1&0\\3&-6&0\end{pmatrix}.$

（Ⅰ）讨论 a, b 为何值时，A 可相似对角化；（Ⅱ）讨论 a, b 为何值时，A 与 B 相似.

\vspace{0.5cm}

\textbf{(22)（本题满分12分）合工大超越十}

设A为3阶实对称矩阵，若存在正交阵 $Q=\frac{1}{\sqrt{2}}\begin{pmatrix}a&1&0\\0&0&c\\1&b&0\end{pmatrix}(b>0,c>0)$，使
\[ \boldsymbol{Q}^{\mathrm{T}}\boldsymbol{A}\boldsymbol{Q}=\boldsymbol{\Lambda}=\begin{pmatrix}{{{2}}}&{{{0}}}&{{{0}}} \\{{{0}}}&{{{0}}}&{{{0}}} \\{{{0}}}&{{{0}}}&{{{1}}}\end{pmatrix}. \]
（Ⅰ）求 a, b, c 及矩阵 A；（Ⅱ）求一个可逆矩阵 P，使 $A = P^{T} P - 2 E$

\vspace{0.5cm}

\textbf{22.(本题满分12分) 张宇八套卷二}

已知矩阵 $A = \begin{pmatrix} 1 & 1 & 1 & 1 \\ 1 & 1 & 1 & 1 \\ 2 & 1 & 0 & -1 \\ -1 & 0 & 1 & 2 \end{pmatrix}$， $B = \begin{pmatrix} 1 & 1 \\ 1 & 1 \\ 0 & 1 \\ 1 & 0 \end{pmatrix}$，矩阵 C 满足 A = BC.

(1) 求矩阵 C;

(2) 计算 $A^{10}$

\vspace{0.5cm}

\textbf{22. 张宇八套卷三}

设 $A=\begin{bmatrix}1&-2&1\\ -2&1&1\\ 1&1&-2\end{bmatrix}$，矩阵 B 满足 AB=A-B，求可逆矩阵 P，使 $P^{-1}(AB)P$ 为对角矩阵，并写出该对角矩阵.

\vspace{0.5cm}

\textbf{22. 张宇八套卷四}

若矩阵 A 的伴随矩阵 $A^{*}=\begin{bmatrix}3&2&-2\\ 0&-1&0\\ a&2&-3\end{bmatrix}$ 相似于矩阵 $B=\begin{bmatrix}-1&0&0\\ -2&1&0\\ 0&0&-1\end{bmatrix}$，其中 $|A|>0$。

(1) 求 a 的值.

(2) 求 $A^{99}$

\vspace{0.5cm}

\textbf{22.(本题满分12分) 张宇八套卷五（同时对角化，难点）}

若可逆线性变换 x = Py 可将二次型 $f(x_{1}, x_{2}) = x_{1}^{2} + 2x_{2}^{2} + 2x_{1}x_{2}$ 化为规范形 $y_{1}^{2} + y_{2}^{2}$，同时将二次型 $g(x_{1}, x_{2}) = -x_{1}^{3} + 2x_{2}^{2} + 2x_{1}x_{2}$ 化为标准形 $k_{1}y_{1}^{2} + k_{2}y_{2}^{2}$，求可逆矩阵 P 及 $k_{1}, k_{2}$ 的值.

\vspace{0.5cm}

\textbf{22 李艳芳三套卷二(谱分解)}

设 A 为 3 阶正交实对称矩阵, $A \neq E$ 且 $|A| = 1, \alpha = (1, 2, 2)^{\mathrm{T}}$ 满足 $A\alpha = \alpha$.

(Ⅰ) 求 A;

（Ⅱ）求可逆线性变换 x = Py，将二次型 $f(x_{1}, x_{2}, x_{3}) = \boldsymbol{x}^{\mathrm{T}} \boldsymbol{A} \boldsymbol{x}$ 化为标准形 $y_{1}^{2} - 4y_{2}^{2} - 16y_{3}^{2}$

\section{第二类 求二次型的最值}

\textbf{22.(本题满分12分) 李林六套卷一}

已知二次型 $f(x_{1},x_{2},x_{3})=x_{1}^{2}+2x_{2}^{2}+ax_{3}^{2}+2x_{1}x_{3}$ 经可逆线性变换 X=PY 化为 $y_{1}^{2}+y_{3}^{2}$.

(Ⅰ) 求 a 的值及可逆矩阵 P;

（Ⅱ）当 $\boldsymbol{X} = (x_{1}, x_{2}, x_{3})^{\mathrm{T}}$ 且 $\boldsymbol{X}^{\mathrm{T}} \boldsymbol{X} = 1$ 时，求 $f(x_{1}, x_{2}, x_{3})$ 的最大值，并求满足 $x_{1} = x_{2} > 0$ 的最大值点.

\vspace{0.5cm}

\textbf{22.(本题满分12分) 张宇八套卷一}

设二次型 $f(x_{1},x_{2},x_{3})=\boldsymbol{x}^{\mathrm{T}}\begin{pmatrix}1&0&6\\4&4&4\\0&8&9\end{pmatrix}\boldsymbol{x}$，其中 $x=\begin{pmatrix}x_{1}\\x_{2}\\x_{3}\end{pmatrix}$

(1) 用正交变换 x = Qy 将其化为标准形，并求出 Q;

(2) 求 $g(x_{1}, x_{2}, x_{3}) = \frac{f(x_{1}, x_{2}, x_{3})}{x_{1}^{2} + x_{2}^{2} + x_{3}^{2}}$ 的最大值，并求出一个最大值点，其中 $x_{1}^{2} + x_{2}^{2} + x_{3}^{2} \neq 0$.

\section{第三类 求 $x^{\top}Ax=0$ 的解}

\textbf{22.(本题满分12分) 李林六套卷二}

设 3 阶实对称矩阵 $\boldsymbol{A}=(\boldsymbol{\alpha}_{1},\boldsymbol{\alpha}_{2},\boldsymbol{\alpha}_{3}),\boldsymbol{r}(\boldsymbol{A})=2$，且满足 $\boldsymbol{\alpha}_{1}+2\boldsymbol{\alpha}_{2}+\boldsymbol{\alpha}_{3}=(3,6,3)^{\mathrm{T}},\boldsymbol{\alpha}_{1}-\boldsymbol{\alpha}_{2}+\boldsymbol{\alpha}_{3}=(-1,1,-1)^{\mathrm{T}}$.

(Ⅰ) 求 A;

（Ⅱ）记 $\boldsymbol{X}=(x_{1},x_{2},x_{3})$，求方程 $\boldsymbol{X}^{\mathrm{T}}(\boldsymbol{A}+\boldsymbol{E})\boldsymbol{X}=0$ 的全部解.

\section{第四类 线性相关无关/线性方程组(或结合特征值特征向量)}

\textbf{22.(本题满分12分) 李林六套卷三}

设 3 阶实对称矩阵 $\boldsymbol{A}=(\boldsymbol{\alpha}_{1},\boldsymbol{\alpha}_{2},\boldsymbol{\alpha}_{3})$ 有二重特征值 $\lambda_{1}=\lambda_{2}=1$，且 $\alpha_{1}+2\alpha_{2}=\alpha_{3}, A^{*}$ 是 A 的伴随矩阵.

（Ⅰ）求一个正交变换 X=QY 化二次型 $f(x_{1},x_{2},x_{3})=\boldsymbol{X}^{\mathrm{T}}\boldsymbol{A}\boldsymbol{X}$ 为标准形.

（Ⅱ）求方程组 $A^{*}X=0$ 的全部解.

\vspace{0.5cm}

\textbf{22.(本题满分12分) 余丙森五套卷一}

已知线性方程组 $\begin{cases}x_{1}+2x_{2}+x_{3}=3,\\2x_{1}+(a+4)x_{2}-5x_{3}=6,\\-x_{1}-2x_{2}+ax_{3}=-3\end{cases}$，有无穷多个解。A为3阶实矩阵，且 $\alpha_{1}=(1,2a,-1)^{\mathrm{T}},\alpha_{2}=(a,a+3,a+2)^{\mathrm{T}},\alpha_{3}=(a-2,-1,a+1)^{\mathrm{T}}$是A的分别对应于特征值 $\lambda_{1}=1,\lambda_{2}=-1,\lambda_{3}=0$的特征向量，

(1) 求矩阵 A;

(2) 求行列式 $\left|A^{2024}+2E\right|$.

\vspace{0.5cm}

\textbf{22.(本题满分12分) 余丙森五套卷二}

设 A 为三阶方阵，并有可逆矩阵 $\boldsymbol{P} = (p_{1}, p_{2}, p_{3}), p_{i} (i = 1, 2, 3)$，为三维列向量，使得
\[ \boldsymbol{P}^{-1}\boldsymbol{A}\boldsymbol{P}=\begin{pmatrix}{{{1}}}&{{{0}}}&{{{0}}} \\{{{0}}}&{{{1}}}&{{{1}}} \\{{{0}}}&{{{0}}}&{{{1}}}\end{pmatrix}. \]
(1) 证明：$p_{1}, p_{2}$ 为 $(E - A)x = 0$ 的解， $p_{3}$ 为 $(E - A)x = -p_{2}$ 的解，且 A 不可相似对角化.

(2) 当 $A = \begin{pmatrix} 2 & -1 & -1 \\ 2 & -1 & -2 \\ -1 & 1 & 2 \end{pmatrix}$ 时，求可逆矩阵 P，使得 $P^{-1} A P = \begin{pmatrix} 1 & 0 & 0 \\ 0 & 1 & 1 \\ 0 & 0 & 1 \end{pmatrix}$.

\vspace{0.5cm}

\textbf{22.(本题满分12分) 余丙森五套卷五}

设有向量组(Ⅰ) $\alpha_{1}=\begin{pmatrix}1\\ 0\\ 1\end{pmatrix},\alpha_{2}=\begin{pmatrix}1\\ 1\\ 2\end{pmatrix},\alpha_{3}=\begin{pmatrix}1\\ 2\\ a\end{pmatrix}$，(Ⅱ) $\beta_{1}=\begin{pmatrix}-1\\ 2\\ 1\end{pmatrix},\beta_{2}=\begin{pmatrix}1\\ 0\\ b\end{pmatrix}$

(1) 问 a, b 为何值时，向量组（Ⅱ）不能由向量组（Ⅰ）线性表示？

(2) 设 $A = \begin{pmatrix} 1 & 1 & 1 \\ 0 & 1 & 2 \\ 1 & 2 & a \end{pmatrix}$, $B = \begin{pmatrix} -1 & 1 \\ 2 & 0 \\ 1 & b \end{pmatrix}$，问 a, b 为何值时，矩阵方程 AX = B 有解，有解时，求出其全部解.

\vspace{0.5cm}

\textbf{22.(本题满分12分) 余丙森五套卷五}

设 3 阶实矩阵 A 满足 $A^{2}-A-2E=O,\left|A\right|=2$

(1) 证明 A 可以对角化；

(2) 如果 A 为实对称矩阵，且 $\boldsymbol{\xi}=(1,1-1)^{\mathrm{T}}$ 是齐次线性方程组 $(A-2\boldsymbol{E})\boldsymbol{x}=\boldsymbol{0}$ 的一个解，求一个对称矩阵 B 使得 $B^{2}=A+E$.

\vspace{0.5cm}

\textbf{(22)（本题满分12分）合工大超越一}

设A为3阶矩阵， $\alpha_{1},\alpha_{2},\alpha_{3}$为3维列向量.已知 $A\alpha_{1}=\alpha_{2},A^{2}\alpha_{1}=\alpha_{3},A^{3}\alpha_{1}=0,\alpha_{3}\neq0$

（Ⅰ）证明：对任意常数 k，向量组 $k\alpha_{1}+\alpha_{2},-\alpha_{2}+\alpha_{3},\alpha_{1}-k\alpha_{3}$ 线性无关；

（Ⅱ）求 A 所有特征值与特征向量，问 A 能否相似于一个对角阵？说明理由；

（Ⅲ）求线性方程组 $A^{T}Ax=A^{T}\alpha_{3}$ 的通解.

\vspace{0.5cm}

\textbf{(22)（本题满分12分）合工大超越八}

已知三阶方阵A的第一行为 $a,b,c,a\neq0,AB=O,B=\begin{pmatrix}1&2&3&1\\2&4&6&2\\3&6&10&4\end{pmatrix}$

$\xi_{1}=\begin{pmatrix}-b\\a\\0\end{pmatrix},\xi_{2}=\begin{pmatrix}-c\\0\\a\end{pmatrix}$，证明：

(Ⅰ) $A\xi_{1}=0$，且B的列向量组与 $\xi_{1},\xi_{2}$等价；

（Ⅱ）求可逆矩阵 P，使 $P^{-1}A^{T}AP=\Lambda.\Lambda$ 为对角阵.

\vspace{0.5cm}

\textbf{22.(本题满分12分) 张宇八套卷八}

已知 3 维列向量 $\xi$ 不是 $A^{2}x = 0$ 的解， $A\xi$ 是 $A^{2}x = 0$ 的解。记 $\mathbf{P} = (\xi, A\xi, A^{2}\xi)$。

(1) 证明 P 可逆;

(2) A 能否相似对角化？若能，求出一个与之相似的对角矩阵，若不能，请说明理由。

\section{第五类 相似、合同的传递性}

\textbf{22.(本题满分12分) 李林六套卷四}

设 $A = \begin{bmatrix} 2 & 1 & 0 \\ 1 & 2 & 0 \\ 0 & 0 & 1 \end{bmatrix}$ 与 $B = \begin{bmatrix} a & b & c \\ 0 & 1 & 0 \\ -1 & -2 & 4 \end{bmatrix}$ 相似.

(Ⅰ) 求 a, b, c 的值；

（Ⅱ）求可逆矩阵 P，使得 $P^{-1}AP=B$

\section{第六类 二次型的正定}
(原文本中此章节标题下无具体题目内容)

\end{document}